% ══════════════════════════════════════════════════════════════
%  Section 9.1: Domain Model --- Blocks, Transactions, and the UTXO Worldview
% ══════════════════════════════════════════════════════════════
\chapter{Chain: Domain Model}
\label{ch:domain-model}

\begin{learnbox}
\begin{itemize}
  \item How blocks and block headers encode the whitepaper's chain of hashes
  \item The UTXO model: why Bitcoin tracks \emph{outputs}, not accounts
  \item How \code{TXInput}, \code{TXOutput}, and \code{Transaction} encode spending
  \item Where ``unspent vs.\ spent'' is tracked (the UTXO set, not the output itself)
\end{itemize}
\end{learnbox}

In this section, we walk through the project's Rust types that represent
Bitcoin's core domain objects. Our goal is to confidently answer ``what
does this type \emph{mean} in whitepaper terms?'' and ``where is it used
in the implementation?''

As we walk through the code, we use these files as our primary reference:

\begin{itemize}
  \item \code{bitcoin/src/primitives/\index{block}block.rs}
  \item \code{bitcoin/src/primitives/\index{transaction}transaction.rs}
  \item \code{bitcoin/src/chain/\index{UTXO}utxo\_set.rs} (for how outputs become ``spendable'')
\end{itemize}

% ──────────────────────────────────────────────────────────────
\section{Whitepaper Term \texorpdfstring{$\to$}{→} Rust Type Map}
\label{sec:wp-rust-map}

\begin{table}[htbp]
\caption{Whitepaper concepts mapped to Rust types}
\label{tab:wp-rust-map}
\centering\small
\begin{tabularx}{0.92\textwidth}{llX}
\toprule
\textbf{Whitepaper Concept} & \textbf{Rust Type / Field} \\
\midrule
Block header          & \code{BlockHeader}                 \\
Previous block hash   & \code{BlockHeader.pre\_block\_hash} \\
Nonce                 & \code{BlockHeader.nonce}            \\
Block hash            & \code{BlockHeader.hash}             \\
Transaction           & \code{Transaction}                  \\
Transaction ID        & \code{Transaction.id}               \\
Input (spend)         & \code{TXInput}                      \\
Output (new coin)     & \code{TXOutput}                     \\
Locking condition     & \code{TXOutput.pub\_key\_hash}       \\
Signature             & \code{TXInput.signature}            \\
\bottomrule
\end{tabularx}
\end{table}

% ──────────────────────────────────────────────────────────────
\section{Domain Object Diagrams}
\label{sec:domain-diagrams}

The diagram below shows the \emph{runtime relationship} between the main
state buckets: the \textbf{\index{blockchain}blockchain} (append-only history), the
\textbf{\index{UTXO}UTXO set} (derived ``what is spendable right now'' state), and
\textbf{\index{transaction}transactions} (which consume old outputs and create new ones).

\begin{figure}[htbp]
\centering
\begin{tikzpicture}[
  node distance = 1.8cm,
  blockchain/.style = {rectangle, rounded corners=6pt, draw=sectioncolor, fill=codebg,
                       text width=6cm, align=center, minimum height=1.2cm, font=\sffamily\small},
  utxo/.style = {rectangle, rounded corners=6pt, draw=brand, fill=brandlight,
                 text width=2.4cm, align=center, minimum height=1cm, font=\sffamily\small},
  arrow/.style = {-{Stealth[length=5pt,width=3.5pt]}, thick, color=brand},
]

% Blockchain
\node[blockchain] (chain) at (0, 3) {
  \textbf{Blockchain} \\[4pt]
  Block 0 (Genesis) \\
  \quad \(\to\) Block 1 \\
  \quad\quad \(\to\) Block 2 \\
  \quad\quad\quad \(\to\) Block N (Tip)
};

% UTXO Set
\node[utxo] (utxo) at (0, 0.5) {
  \textbf{UTXO Set} \\[4pt]
  \texttt{(txid, vout)} \\
  \quad \(\to\) TXOutput \\
  \texttt{(txid, vout)} \\
  \quad \(\to\) TXOutput
};

% Connecting arrow and label
\draw[arrow] (chain.south) -- (utxo.north);
\node[font=\small\sffamily\itshape, color=sectioncolor, align=center, right=0.3cm of chain.south] {
  maintains \\ derived state
};

\end{tikzpicture}
\caption{Blockchain Domain Model: Chain History and Derived UTXO State}
\label{fig:ch09-1-domain-model}
\end{figure}

% ──────────────────────────────────────────────────────────────
\section{Key Code Reading Checklist}
\label{sec:reading-checklist}

As we follow the Rust code, we should be able to answer these questions
from the structs alone:

\begin{itemize}
  \item Where does the txid come from? (Search for \code{Transaction::hash()}.)
  \item Where does the signature live? (Look at \code{TXInput.signature}.)
  \item What does an output lock to? (\code{TXOutput.pub\_key\_hash}.)
  \item Where is ``unspent vs.\ spent'' tracked? (Not on the output---in \code{UTXOSet}.)
\end{itemize}

% ──────────────────────────────────────────────────────────────
\section{Step-by-Step Code Walkthrough}
\label{sec:walkthrough}

\textbf{Goal}: build a correct mental model of what the core structs
\emph{mean}, and how the bytes flow between modules.

\begin{infobox}{Source}
\code{bitcoin/src/primitives/block.rs},
\code{bitcoin/src/primitives/transaction.rs},
\code{bitcoin/src/chain/utxo\_set.rs}
\end{infobox}

% ── Step 1: Block and BlockHeader ────────────────────────────
\subsection{Step 1: Block and BlockHeader}
\label{sec:step1-block}

In the Bitcoin whitepaper, a block is a \textbf{timestamped commitment}
to two things: the previous block (so the history is tamper-evident by
hash chaining) and a set of transactions (so the block anchors ownership
transfers into the history).

In our implementation, \code{BlockHeader} is the smallest header we need
to make those ideas concrete in code. In practice, the header matters
because it defines the \textbf{proof-of-work input}: mining repeatedly
hashes a byte sequence derived from the header (plus a transaction
commitment) until it finds a valid \code{nonce}.

Those bytes are assembled from the header fields in
\code{ProofOfWork::prepare\_data(...)}, which is why \code{BlockHeader}
contains exactly the fields it does. We cover the proof-of-work loop
and the exact hashed byte layout in Chapter~15 (Block Lifecycle and
Mining).

\begin{listing}[htbp]
\caption{Block and BlockHeader struct layout (\code{block.rs})}
\label{lst:9-1}
\begin{minted}[fontsize=\footnotesize]{rust}
pub struct BlockHeader {
    timestamp: i64,
    pre_block_hash: String,
    hash: String,
    nonce: i64,
    height: usize,
}

pub struct Block {
    header: BlockHeader,
    transactions: Vec<Transaction>,
}
\end{minted}
\end{listing}

What the fields mean:
\begin{itemize}
  \item \code{pre\_block\_hash}: link to the previous block (whitepaper \S3/\S4).
  \item \code{nonce} + \code{hash}: PoW search result (whitepaper \S4).
  \item \code{timestamp}: ordering signal; also part of the PoW header bytes.
  \item \code{height}: convenience for iteration/selection; not a Bitcoin header field.
\end{itemize}

\begin{notebox}[Project-specific note]
In our implementation we do not use a \textbf{Merkle root} field. We
instead concatenate transaction IDs and hash once using
\code{Block::hash\_transactions()}. This keeps the code readable but
does not support Merkle proofs or SPV.
\end{notebox}

Creating a new block: in our implementation,
\code{Block::new\_block(pre\_block\_hash, transactions, height)} is the
entry point that \emph{defines} the header boundary. It takes the parent
link and the transaction list as inputs, sets derived header fields such
as \code{timestamp}, and then finalizes \code{nonce} and \code{hash} by
running PoW.

\begin{listing}[htbp]
\caption{\code{Block::new\_block} --- creating and mining a block}
\label{lst:9-1-2}
\begin{minted}[fontsize=\footnotesize]{rust}
pub fn new_block(
    pre_block_hash: String,
    transactions: &[Transaction],
    height: usize
) -> Block {
    let header = BlockHeader {
        timestamp: crate::current_timestamp(),
        pre_block_hash,
        hash: String::new(),
        nonce: 0,
        height,
    };
    let mut block = Block {
        header,
        transactions: transactions.to_vec(),
    };
    let pow = ProofOfWork::new_proof_of_work(block.clone());
    let (nonce, hash) = pow.run();
    block.header.nonce = nonce;
    block.header.hash = hash;
    block
}
\end{minted}
\end{listing}

\begin{checkpointbox}
\code{pre\_block\_hash} is the chain link. \code{nonce} and \code{hash}
are filled by PoW (\code{ProofOfWork::run()}), not supplied by the
caller. \code{height} lives in the header as a convenience (not a
Bitcoin header field).
\end{checkpointbox}

% ── Step 2: Transaction, Inputs, Outputs ─────────────────────
\subsection{Step 2: Transaction, Inputs, and Outputs}
\label{sec:step2-transaction}

Bitcoin is not an ``account system.'' It is a system of \textbf{outputs},
and spending means consuming previous outputs as inputs.

\begin{infobox}{Methods involved}
\code{Transaction::new\_utxo\_transaction()},
\code{Transaction::hash()},
\code{Transaction::sign()}
\end{infobox}

In this step, we do two things: (1)~define what a transaction \emph{is}
by reading the struct fields, and (2)~show how a spending transaction is
\emph{constructed} in code.

\begin{listing}[htbp]
\caption{Transaction, TXInput, and TXOutput structs (\code{transaction.rs})}
\label{lst:9-2-structs}
\begin{minted}[fontsize=\footnotesize]{rust}
#[derive(Clone, Default, Serialize, Deserialize)]
pub struct TXInput {
    txid: Vec<u8>,
    vout: usize,
    signature: Vec<u8>,
    pub_key: Vec<u8>,
}

#[derive(Clone, Serialize, Deserialize)]
pub struct TXOutput {
    value: i32,
    in_global_mem_pool: bool,
    pub_key_hash: Vec<u8>,
}

#[derive(Clone, Default, Serialize, Deserialize)]
pub struct Transaction {
    id: Vec<u8>,
    vin: Vec<TXInput>,
    vout: Vec<TXOutput>,
}
\end{minted}
\end{listing}

They reveal the project's UTXO worldview: a transaction spends prior
outputs by referencing an outpoint $(txid, vout)$ in \code{TXInput}, and
it creates new spendable outputs in \code{TXOutput}.

\begin{itemize}
  \item A spend is identified by an outpoint $(txid, vout)$: \code{TXInput.\{txid, vout\}}.
  \item Authorization is per-input: \code{TXInput.signature} + \code{TXInput.pub\_key}.
  \item Outputs are locked by a pubkey-hash: \code{TXOutput.pub\_key\_hash} (no script engine).
\end{itemize}

How a spending transaction is constructed:
\begin{enumerate}
  \item \textbf{Select spendable outputs}: query the UTXO set for outpoints locked to the sender.
  \item \textbf{Create inputs}: for each selected outpoint, create a \code{TXInput}.
  \item \textbf{Create outputs}: payment to recipient, plus optional change back to sender.
  \item \textbf{Compute the txid}: hash a serialized copy with \code{id = []} (non-circular).
  \item \textbf{Sign}: produce per-input signatures and store into \code{TXInput.signature}.
\end{enumerate}

The concrete entry point is \code{Transaction::new\_utxo\_transaction(...)}
(full implementation in Chapter~14: Transaction Lifecycle).

% ── Step 3: UTXO Set ────────────────────────────────────────
\subsection{Step 3: History vs.\ Derived State (the UTXO Set)}
\label{sec:step3-utxo}

\textbf{UTXO (Unspent Transaction Output)}: a specific transaction
output that is currently spendable---an output identified by an outpoint
$(txid, vout)$ that has \textbf{not} been consumed by any later
transaction input.

The \textbf{blockchain} is the append-only history (blocks and the
transactions they contain). The \textbf{UTXO set} is derived state: a
compact, query-friendly view of ``what is spendable right now?'' computed
from that history.

We keep the UTXO set separate so the node can answer spendability
questions efficiently without rescanning the entire chain every time it
validates a transaction or connects a block.

\begin{listing}[htbp]
\caption{UTXOSet struct (\code{utxo\_set.rs})}
\label{lst:9-3-utxo}
\begin{minted}[fontsize=\footnotesize]{rust}
const UTXO_TREE: &str = "chainstate";

pub struct UTXOSet {
    blockchain: BlockchainService,
}

impl UTXOSet {
    pub fn new(blockchain: BlockchainService) -> UTXOSet {
        UTXOSet { blockchain }
    }
}
\end{minted}
\end{listing}

\begin{notebox}
Spendability is not a field on \code{TXOutput}; it is tracked by
membership in the UTXO database. This is the same design Bitcoin Core
uses: the \code{chainstate} LevelDB stores only unspent outputs.
\end{notebox}

\begin{refbox}{See the full implementation}
The complete \code{find\_spendable\_outputs()} method, UTXO rebuild
logic, and rollback handling are covered in Chapter~13 (UTXO Set).
\end{refbox}
